\subsection{Admissibility criteria and scope of the guarantees}
\label{sec:requirements}
A small prediction error is not sufficient to define an admissible hyperelastic surrogate. We organize the construction around six criteria, following the distinctions made in constitutive PANN formulations~\cite{klein2022,linden2023}. The enumeration is an organizing device for this paper, not a universal checklist or a claim that polyconvexity is necessary for every physical material:
\begin{enumerate}
 \item[(C1)] \emph{Potential structure:} stress is the derivative of a single-valued stored energy in work-conjugate variables.
 \item[(C2)] \emph{Objectivity:} a superposed rigid rotation does not change that energy.
 \item[(C3)] \emph{Stress symmetry and anisotropy:} the material stress is symmetric, while directional dependence is allowed without imposing an unverified material symmetry group.
 \item[(C4)] \emph{Reference normalization:} the undeformed state has zero energy and zero stress.
 \item[(C5)] \emph{Polyconvexity:} the energy admits a convex representation in deformation minors, providing a sufficient rank-one convexity condition.
 \item[(C6)] \emph{Growth:} energy diverges under volume collapse, unbounded volume expansion, and unbounded distortion at bounded positive volume.
\end{enumerate}
The ICNN and ICKAN constructions below satisfy these criteria for every finite parameter set obeying their architectural constraints. Global nonnegative energy is an additional condition checked separately for the selected parameters. Constitutive identities, sufficient bounds, and sampled numerical checks consequently play different roles.

\paragraph{Potential structure (C1).}
For a differentiable hyperelastic energy, stress is its derivative in work-conjugate variables. In the convention of Eq.~\eqref{eq:voigt},
\begin{equation}
 \ssv(\ee)=\nabla_{\ee}W(\ee),\qquad
 \oint_{\Gamma}\ssv\cdot\dd\ee=0 .
 \label{eq:conservative_requirement}
\end{equation}
For the isothermal purely elastic setting considered here, (C1) also gives $\ssv\cdot\dot{\ee}-\dot W=0$: the mechanical dissipation vanishes. The closed-path identity is exact for a single-valued differentiable potential. On a simply connected strain domain, a continuously differentiable stress map has such a potential if its Jacobian is symmetric. Direct stress regression need not impose the cross-derivative identities
$\partial s_i/\partial e_j=\partial s_j/\partial e_i$~\cite{asad2022}. A small pointwise stress error alone does not establish them.

\paragraph{Objectivity and the three meanings of symmetry (C2)--(C3).}
For a superposed proper rotation $\bm Q$, objectivity requires
\begin{equation}
 W(\bm Q\F)=W(\F),\qquad
 \bm P(\bm Q\F)=\bm Q\bm P(\F).
\end{equation}
Using $\C=\F^T\F$ as an input to an energy enforces the first identity; differentiating the energy then gives the stress transformation. A direct map from $\E$ to the material stress $\Sg$ can also be objective without possessing a potential. Objectivity therefore does not distinguish Regression from Free in this study. The symmetry $\Sg=\Sg^T$ follows from differentiation with respect to symmetric strain for the energy models and is built into the three-component output of Regression. It implies a symmetric Cauchy stress $\bm\sigma=J^{-1}\F\Sg\F^T$, as required by angular-momentum balance in the present continuum.

Neither stress symmetry nor objectivity imposes a material symmetry group. A right action $\F\mapsto\F\bm Q_m$ leaves the energy unchanged only when $\bm Q_m$ belongs to that group. Directional features are fixed in the reference material and are unchanged by a superposed spatial rotation. Finally, the constitutive tangent $\bm D=\partial\ssv/\partial\ee\in\R^{3\times3}$ measures how the three stress components change with engineering strain, with entries $D_{ij}=\partial s_i/\partial e_j$. For an energy-based law it is the Hessian $\partial^2_{\ee\ee}W$. Its \emph{major symmetry}, $\bm D=\bm D^T$, is a third property: it follows from a twice differentiable scalar energy, not merely from a symmetric stress tensor.

\paragraph{Reference normalization (C4).}
Reference normalization imposes
\begin{equation}
 W(\I)=0,\qquad \bm P(\I)=\bm0.
\end{equation}
The first condition fixes an arbitrary additive energy constant. The second is a substantive stress condition and depends on how the energy is corrected. Subtracting a linear function of $\E$ preserves a potential but need not preserve polyconvexity in $\F$. An affine correction in independent deformation minors, in contrast, preserves convexity in those variables. This is the representation introduced next and used by our normalization.

\paragraph{Polyconvexity and finite-strain stability (C5).}
Potential structure does not impose a stability condition. The Legendre--Hadamard condition for a twice differentiable energy is
\begin{equation}
 \partial^2_{\F\F}W(\F)[\bm a\otimes\bm b,\bm a\otimes\bm b]\geq0
 \quad\hbox{for all }\bm a,\bm b .
 \label{eq:lh_background}
\end{equation}
Here $\bm a,\bm b\in\R^2$ are arbitrary vectors, $(\bm a\otimes\bm b)_{ij}=a_i b_j$ is their outer product, and $t\in\R$ parametrizes the perturbation. The condition tests the curvature of the scalar function $t\mapsto W(\F+t\bm a\otimes\bm b)$ at $t=0$: the energy must not curve downward along a rank-one perturbation. Strict positivity for nonzero directions is a stronger statement than this nonnegative condition. In two dimensions, a sufficient condition is a representation
\begin{equation}
 W(\F)=\mathcal P(\F,\det\F),\qquad
 \mathcal P\text{ convex in its independent arguments}.
 \label{eq:polyconvex_definition}
\end{equation}
Here $\mathcal P$ is defined on the convex set $\R^{2\times2}\times(0,\infty)$, treating its matrix and scalar arguments as independent; only its evaluation in $W$ sets the second argument equal to $\det\F$. This distinction is essential: we do not require $W$ to be convex in $\F$ itself. The cofactor is linear in the entries of a $2\times2$ matrix, so it need not be introduced as an additional independent argument. Along a rank-one perturbation, both $\F$ and $\det\F$ vary affinely with the perturbation parameter. The resulting path is therefore a straight line in the arguments of $\mathcal P$, whose convexity gives Eq.~\eqref{eq:lh_background}. Appendix~\ref{app:finite_strain_curvature} supplies the derivation. Testing finitely many rank-one directions cannot establish polyconvexity.

The standard three-dimensional definition uses $(\F,\cof\F,\det\F)$~\cite{ball1977,Ciarlet1988}. The guarantee established here is two-dimensional; it does not establish an admissible three-dimensional extension with arbitrary out-of-plane deformation.

Convexity in $\E$ is a different property because $\E=\tfrac12(\F^T\F-\I)$ depends quadratically on $\F$. Consequently, the energy curvature along a linear perturbation of $\F$ includes both a strain-tangent contribution and a contribution from the curvature of the strain path, weighted by the current stress. The latter can be negative in compression, so a positive-semidefinite strain Hessian $\bm D$ alone does not guarantee Eq.~\eqref{eq:lh_background}~\cite{asad2022,GaoNeffRoventaThiel2017}. Convexity in $\C$ and in $\E$ are equivalent up to their affine change of variables, but neither implies polyconvexity. Appendix~\ref{app:finite_strain_curvature} derives this stress-dependent curvature and gives an analytical counterexample.

\paragraph{Growth and the limits of the guarantees (C6).}
Finally, the inequalities $W(\F)\geq0$, the barrier $W\to\infty$ as $J\to0^+$, and the existence of minimizers are not synonyms. Ciarlet's presentation of Ball's existence argument~\cite{Ciarlet1988}, especially Sections 4.9 and 7.7, includes coercivity, admissible boundary data, and finite-energy competitors in addition to polyconvexity. Neither the existence theory nor a constitutive barrier guarantees convergence of an unconstrained Newton step. In this work, architectural statements, saved-parameter bounds, and numerical observations are therefore identified separately.
